\documentclass[12pt]{article}
\usepackage[utf8]{inputenc}
\usepackage[T1]{fontenc}
\usepackage{hyperref}
\title{Xet\'a Language Resources for K-12 Curriculum Development}
\author{Compiled by Claude Sonnet 5.5 (claude-sonnet-5-5) for Andr\'as Kornai}
\date{}
\begin{document}
\maketitle

\section{Introduction}
Xet\'a (ISO 639-3 xet; Glottocode xeta1241) is a dormant Tupi-Guarani language of northwestern Paran\'a, Brazil. Sources found in search place it near the Guarani group, especially M\-by\'a, in phonology and lexicon. Traditional territory lay between the Ivaí River and its tributaries. Only a handful of children survived the contact period, and one source gives about 69 people identifying as Xet\'a in 2020. The search results did not turn up a published dictionary, a textbook or any school materials.

\section{Primary Educational Resources}
No primer, textbook or classroom materials in Xet\'a were found. A revitalization effort is described under Community Resources below.

\section{Linguistic and Reference Materials}
\begin{itemize}
\item Wikipedia overview of the language: \url{https://en.wikipedia.org/wiki/Xeta_language}
\item Povos Ind\'igenas no Brasil (Instituto Socioambiental) entry on the Xet\'a people: \url{https://pib.socioambiental.org/en/Povo:Xet\'a}
\item Indios.org.br entry on the Xet\'a people: \url{https://indios.org.br/en/povo/xeta}
\item A repository item at eduCAPES (UFPR) relating to Xet\'a: \url{https://educapes.capes.gov.br/handle/1884/37035?mode=full}. I did not confirm its content.
\end{itemize}
I did not find a dictionary or wordlist. Historical wordlists probably exist in Portuguese-language ethnographic and linguistic literature, but I saw no URL for them.

\section{Religious and Community Resources}
A conference paper from UCLA describes ``Jan\'e Reko Paranuh\'a'', a teacher-education and language-revival programme run through the Universidade Estadual de Maring\'a: \url{https://conferences.gseis.ucla.edu/simposio/ponencias-presentations/faustinorosangela-celia-keros-gustavo-mileski-universidade-estadual-de-maringa-brasil-.-jane-reko-paranuha-a-experiencia-com-a-revitalizacao-da-lingua-xeta-no-parana-jane-reko-paranuha-an-experience-in-indigenous-teachers-formation-and-the-re/jane-reko-paranuha-an-experience-in-indigenous-teachers-formation-and-the-revival-of-the-xeta-language-in-the-state-of-parana-brazil/at_download/file}

\section{Digital and Audio Resources}
No audio, video or software resources were found.

\section{Curriculum Development Context}
Little documentation exists, so a curriculum would have to be built from historical wordlists and the related Guarani languages, working with the Xet\'a community and the Maring\'a programme. Any such work should be community-led. Two anthropological and historical papers also turned up: \url{https://www.anchietano.unisinos.br/publicacoes/antropologia/volumes/073/010.pdf} and \url{https://periodicos.pucpr.br/ceb/article/download/22509/21593}.
\end{document}